\documentclass[10pt,a4paper]{amsart}
\usepackage{amsmath,amssymb,amsthm,graphicx}
\usepackage{verbatim}
\newenvironment{reference}{\comment}{\endcomment}
%\newenvironment{reference}{}{}
\newenvironment{slogan}{\comment}{\endcomment}
\newenvironment{history}{\comment}{\endcomment}

% For commutative diagrams we use Xy-pic
\usepackage[all]{xy}

% We use 2cell for 2-commutative diagrams.
\xyoption{2cell}
% Theorem environments.
%
\theoremstyle{plain}
\newtheorem{theorem}[subsection]{Theorem}
\newtheorem{proposition}[subsection]{Proposition}
\newtheorem{lemma}[subsection]{Lemma}

\theoremstyle{definition}
\newtheorem{definition}[subsection]{Definition}
\newtheorem{example}[subsection]{Example}
\newtheorem{exercise}[subsection]{Exercise}
\newtheorem{situation}[subsection]{Situation}

\theoremstyle{remark}
\newtheorem{remark}[subsection]{Remark}
\newtheorem{remarks}[subsection]{Remarks}

\numberwithin{equation}{subsection}

% Macros
%
\def\lim{\mathop{\mathrm{lim}}\nolimits}
\def\colim{\mathop{\mathrm{colim}}\nolimits}
\def\Spec{\mathop{\mathrm{Spec}}}
\def\Hom{\mathop{\mathrm{Hom}}\nolimits}
\def\Ext{\mathop{\mathrm{Ext}}\nolimits}
\def\SheafHom{\mathop{\mathcal{H}\!\mathit{om}}\nolimits}
\def\SheafExt{\mathop{\mathcal{E}\!\mathit{xt}}\nolimits}
\def\Sch{\mathit{Sch}}
\def\Mor{\mathop{\mathrm{Mor}}\nolimits}
\def\Ob{\mathop{\mathrm{Ob}}\nolimits}
\def\Sh{\mathop{\mathit{Sh}}\nolimits}
\def\NL{\mathop{N\!L}\nolimits}
\def\CH{\mathop{\mathrm{CH}}\nolimits}
\def\proetale{{pro\text{-}\acute{e}tale}}
\def\etale{{\acute{e}tale}}
\def\QCoh{\mathit{QCoh}}
\def\Ker{\mathop{\mathrm{Ker}}}
\def\Im{\mathop{\mathrm{Im}}}
\def\Coker{\mathop{\mathrm{Coker}}}
\def\Coim{\mathop{\mathrm{Coim}}}

% Boxtimes
%
\DeclareMathSymbol{\boxtimes}{\mathbin}{AMSa}{"02}

%
% Macros for moduli stacks/spaces
%
\def\QCohstack{\mathcal{QC}\!\mathit{oh}}
\def\Cohstack{\mathcal{C}\!\mathit{oh}}
\def\Spacesstack{\mathcal{S}\!\mathit{paces}}
\def\Quotfunctor{\mathrm{Quot}}
\def\Hilbfunctor{\mathrm{Hilb}}
\def\Curvesstack{\mathcal{C}\!\mathit{urves}}
\def\Polarizedstack{\mathcal{P}\!\mathit{olarized}}
\def\Complexesstack{\mathcal{C}\!\mathit{omplexes}}
% \Pic is the operator that assigns to X its picard group, usage \Pic(X)
% \Picardstack_{X/B} denotes the Picard stack of X over B
% \Picardfunctor_{X/B} denotes the Picard functor of X over B
\def\Pic{\mathop{\mathrm{Pic}}\nolimits}
\def\Picardstack{\mathcal{P}\!\mathit{ic}}
\def\Picardfunctor{\mathrm{Pic}}
\def\Deformationcategory{\mathcal{D}\!\mathit{ef}}
\usepackage{cleveref}
\crefname{originallemma}{original lemma}{original lemmas}
\Crefname{originallemma}{Original lemma}{Original lemmas}
\makeatletter
\newlabel{descent-lemma-open-fpqc-covering}{{03N0}{source}{}{sourcepointer-1}{}}
\newlabel{descent-lemma-open-fpqc-covering@cref}{{[originallemma][0][]03N0}{[1][0][]source}{}{}{}}
\makeatother
\crefname{originalreference}{original reference}{original references}
\Crefname{originalreference}{Original reference}{Original references}
\makeatletter
\newlabel{lemma-compare-cohomology-nonabelian}{{099X}{source}{}{sourcepointer-2}{}}
\newlabel{lemma-compare-cohomology-nonabelian@cref}{{[originalreference][0][]099X}{[1][0][]source}{}{}{}}
\makeatother
\crefname{originallemma}{original lemma}{original lemmas}
\Crefname{originallemma}{Original lemma}{Original lemmas}
\makeatletter
\newlabel{sites-modules-lemma-locally-constant}{{093T}{source}{}{sourcepointer-3}{}}
\newlabel{sites-modules-lemma-locally-constant@cref}{{[originallemma][0][]093T}{[1][0][]source}{}{}{}}
\makeatother
\title{Pullback along the comparison from the pro-etale to the etale site is an equivalence on locally constant sheaves of finite presentation}
\author{The Stacks Project Authors}
\date{stacks-project:proetale.tex:local-snapshot-20261012; GNU FDL 1.2 or later}
\begin{document}
\maketitle
\noindent\textit{Editorial note.} All delivered mathematical text is the author's own, copied from the pinned Stacks Project chapter file proetale.tex; no statement or proof step was rewritten, reordered or removed.\par\medskip
\noindent\textit{Editorial note.} The counted proof is the author's complete argument: the splitting of the base into disjoint open pieces on which the trivializing module is constant, the passage to the torsor of isomorphisms with a constant sheaf, and the conclusion from the agreement of torsors between the etale and pro-etale topologies.\par\medskip
\noindent\textit{Editorial note.} The label lemma-compare-locally-constant names the counted statement, and the proof environment immediately following it in the source is delivered verbatim as its proof.\par\medskip
\noindent\textit{Editorial note.} The counted statement is carried by the author's own lemma environment in the chapter source.\par\medskip
\noindent\textit{Editorial note.} Pre-open gate: python operations/agent-authoring/tools/precard\_gate.py --source proetale.tex --label lemma-compare-locally-constant --start 4337 --end 4388 --local-id 3588 returned RESULT=OK (L3588 free, occupancy FREE, public interval clear). Candidate taken from sources/stacks-proetale/candidates-halfA.json.\par\medskip
\noindent\textit{Editorial note.} The statement and proof are copied from the pinned chapter file; the cited inputs are kept as external pointers with their canonical Stacks tags.\par\medskip
\noindent\textit{Allowed standard prerequisites (editorial).}\begin{itemize}
\item Openness of the image of an fpqc covering. Descent, the lemma that the image of an fpqc covering morphism is open. Supplies the openness of the images of the pieces of a pro-etale covering, used to cut X into the disjoint open pieces U\_M.
\item The sheaf of isomorphisms with a constant sheaf is constant. Modules on Sites, the lemma identifying the sheaf of isomorphisms between two locally constant sheaves with a constant sheaf. Supplies the identification of the sheaf G of automorphisms of the constant sheaf M with the constant sheaf on the group of module automorphisms of M.
\item Agreement of torsors for the etale and pro-etale topologies. Same chapter, the lemma comparing the first cohomology with nonabelian coefficients for the pro-etale and the etale topology. Supplies the fact that the torsor of isomorphisms between the constant sheaf and F has sections etale locally.
\end{itemize}\medskip
\ifdefined\hypertarget\hypertarget{sourcepointer-1}{}\fi
\noindent\textit{Source reference.} Original lemma 03N0: sources/stacks-chow/tags.txt, source lines 4768--4768. Prerequisite pointer invoked at proof line(s) 4369; it names a standard fact of a sibling chapter that the counted argument uses, and every step of the counted proof built on it is retained verbatim in the delivered body.\par\medskip
\ifdefined\hypertarget\hypertarget{sourcepointer-2}{}\fi
\noindent\textit{Source reference.} Original reference 099X: sources/stacks-proetale/pinned-original-source.tex, source lines 4226--4226. Prerequisite pointer invoked at proof line(s) 4384; it names a standard fact of the same chapter that the counted argument uses, and every step of the counted proof built on it is retained verbatim in the delivered body.\par\medskip
\ifdefined\hypertarget\hypertarget{sourcepointer-3}{}\fi
\noindent\textit{Source reference.} Original lemma 093T: sources/stacks-chow/tags.txt, source lines 11526--11526. Prerequisite pointer invoked at proof line(s) 4379; it names a standard fact of a sibling chapter that the counted argument uses, and every step of the counted proof built on it is retained verbatim in the delivered body.\par\medskip
\begin{lemma}
\label{lemma-compare-locally-constant}
Let $X$ be a scheme. Let $\Lambda$ be a ring.
The functor $\epsilon^{-1}$ defines an equivalence of categories
$$
\left\{
\begin{matrix}
\text{locally constant sheaves}\\
\text{of }\Lambda\text{-modules on }X_\etale\\
\text{of finite presentation}
\end{matrix}
\right\}
\longleftrightarrow
\left\{
\begin{matrix}
\text{locally constant sheaves}\\
\text{of }\Lambda\text{-modules on }X_\proetale\\
\text{of finite presentation}
\end{matrix}
\right\}
$$
\end{lemma}

\begin{proof}
Let $\mathcal{F}$ be a locally constant sheaf of $\Lambda$-modules
on $X_\proetale$ of finite presentation. Choose a pro-\'etale covering
$\{U_i \to X\}$ such that $\mathcal{F}|_{U_i}$ is constant, say
$\mathcal{F}|_{U_i} \cong \underline{M_i}_{U_i}$.
Observe that $U_i \times_X U_j$ is empty if $M_i$ is not isomorphic
to $M_j$.
For each $\Lambda$-module $M$ let $I_M = \{i \in I \mid M_i \cong M\}$.
As pro-\'etale coverings are fpqc coverings and by
Descent, Lemma \ref{descent-lemma-open-fpqc-covering}
we see that $U_M = \bigcup_{i \in I_M} \Im(U_i \to X)$
is an open subset of $X$. Then $X = \coprod U_M$ is a disjoint
open covering of $X$. We may replace $X$ by $U_M$ for some $M$ and
assume that $M_i = M$ for all $i$.

\medskip\noindent
Consider the sheaf $\mathcal{I} = \mathit{Isom}(\underline{M}, \mathcal{F})$.
This sheaf is a torsor for
$\mathcal{G} = \mathit{Isom}(\underline{M}, \underline{M})$.
By Modules on Sites, Lemma \ref{sites-modules-lemma-locally-constant}
we have $\mathcal{G} = \underline{G}$
where $G = \mathit{Isom}_\Lambda(M, M)$.
Since torsors for the \'etale topology
and the pro-\'etale topology agree by
Lemma \ref{lemma-compare-cohomology-nonabelian}
it follows that $\mathcal{I}$ has sections \'etale locally on $X$.
Thus $\mathcal{F}$ is \'etale locally a constant sheaf which is
what we had to show.
\end{proof}

\end{document}
